\chapter{APIs}
\label{ch:apis}

This chapter presents the public interface in terms of the objects and rules
of \crefrange{ch:definitions}{ch:semantics}: for each entry point, a
signature, a contract, and which $\Lom$ rules the call fires.

The contracts use a fixed vocabulary:

\begin{description}[leftmargin=1.9cm,style=nextline,font=\normalfont\bfseries]
  \item[Requires] what must hold on entry. Split by \emph{enforcement}:
    \textsc{chk} is verified and reported through the return value,
    \textsc{ast} is verified and aborts on failure, and \textsc{res} is an
    unverified program restriction from \cref{sec:l0-restrictions}.
  \item[Ensures] the postcondition, in terms of $\StoreH$ or the
    configuration.
  \item[Frame] the components of \cref{def:configuration} the call may
    modify.
  \item[Blocks] whether the call may wait, and on what.
\end{description}

The \textsc{res} rows are the ones to read carefully. They are the price of
the memory model of \cref{sec:memory-model}, and none is currently checked
\devmark{dev:lookahead-uncheck}.

% ===========================================================================
\section{Session lifecycle}
\label{sec:api-session}

\begin{lstlisting}
int  dstm_init(int *argc, char **argv[]);
int  dstm_init_comm(int *argc, char **argv[], MPI_Comm comm);
int  dstm_start(dstm_distance_t *dist);
int  dstm_start_custom(dstm_distance_t *dist, dstm_opts_t opts);
int  dstm_fini(dstm_time_t last_time);
int  dstm_barrier(dstm_time_t last_time);
int  dstm_finalize(void);
\end{lstlisting}

\code{dstm\_init} establishes the process-wide communicator but begins no
\gls{session}; \code{dstm\_start} begins one, fixing $\cgraph$ and allocating
every region. Before the first \code{dstm\_start} there is no region, so no
versioned memory exists.

\begin{table}[htbp]
\centering\small
\begin{tabularx}{\textwidth}{@{}l l >{\raggedright\arraybackslash}X@{}}
\toprule
\multicolumn{3}{@{}l}{\code{dstm\_start(dist)} --- open a session} \\
\midrule
Requires & \textsc{chk} & \code{dist} non-\code{NULL}, of size $\worldsize$;
  region size $\ge$ 64 bytes \\
         & \textsc{res} & $\Dmat$ encodes a well-formed \gls{comm-graph}
  (\cref{def:comm-graph}). Only the diagonal and $\Dmat_{ij} \ge -1$ are
  checked \devmark{dev:zero-distance} \devmark{dev:asymmetric} \\
         & \textsc{res} & collective: all processes call it, in the same order
  (\cref{res:collective}) \\
Ensures  & & $\clk{p} = 0$, $\kn{p}{q} = 0$, $\Vers{p} = \emptyset$,
  $\Pend{p} = \emptyset$, $\Net = \emptyset$; i.e.\ $\Cfg_0$ of
  \cref{def:configuration} \\
Frame    & & the whole configuration \\
Blocks   & & no \\
\bottomrule
\end{tabularx}
\end{table}

\begin{table}[htbp]
\centering\small
\begin{tabularx}{\textwidth}{@{}l l >{\raggedright\arraybackslash}X@{}}
\toprule
\multicolumn{3}{@{}l}{\code{dstm\_barrier(last\_time)} --- establish quiescence} \\
\midrule
Requires & \textsc{res} & $\ell \ge \max_{o} \opeff{o}$ over \emph{all} processes,
  and identical on every process (\cref{res:terminate})
  \devmark{dev:clock-conflation} \\
         & \textsc{res} & collective (\cref{res:collective}) \\
Ensures  & & quiescence (\cref{def:quiescence}): $\Pend{p} = \emptyset$ at
  every process and no data message in $\Net$
  (\cref{thm:barrier-correct}) \\
Frame    & & $\Pend{p}$, $\Vers{p}$, $\Net$, $\mathit{drained}_p$ \\
Blocks   & & yes, until the three-round ring completes
  (\rn{Bar-One}, \rn{Bar-Two}, \rn{Bar-Exit}) \\
\bottomrule
\end{tabularx}
\end{table}

Passing too small an $\ell$ is the commonest way to hang a DSTM program: the
second round relies on $\mathit{drained}_p$ to let the final \glspl{version}
publish, so a process still expecting an \gls{effective-time} beyond $\ell$
never resolves its pending get. \code{dstm\_fini} is \code{dstm\_barrier} plus
teardown of the region, graph and \gls{session} communicator, and inherits the
restriction.

% ===========================================================================
\section{Distance matrices}
\label{sec:api-distance}

\begin{lstlisting}
int dstm_distance_create_ring(dstm_distance_t **ret, size_t size);
int dstm_distance_create_complete(dstm_distance_t **ret, size_t size);
int dstm_distance_create_torus(dstm_distance_t **ret, size_t size);
/* ... star, line, mesh, tree, hypercube, dragonfly ... */
int dstm_distance_create_from_matrix(dstm_distance_t **ret, size_t size,
                                     const int *matrix);
int dstm_distance_get(dstm_distance_t *dst, int i, int j);
void dstm_distance_destroy(dstm_distance_t **dist);
\end{lstlisting}

\alsoterm{distance-matrix}%
A \code{dstm\_distance\_t} is the API-level encoding of a
\gls{comm-graph}: $\Dmat_{ij} = \lag{i}{j}$ when a channel joins the two
processes, $\Dmat_{ii} = 0$, and $-1$ where there is no channel. Note that
$-1$ encodes \emph{absence of a channel}, not an infinite lag --- the two
ranks may well be joined by a path
(\cref{rem:path-weight-metric}).

The named constructors assign equal weight to every channel, so they are
so (G1)--(G3) hold by construction.
\code{dstm\_distance\_create\_from\_matrix} is the escape hatch for
application-derived topologies and is the only place where a malformed graph
can enter the system.

\begin{table}[htbp]
\centering\small
\begin{tabularx}{\textwidth}{@{}l l >{\raggedright\arraybackslash}X@{}}
\toprule
\multicolumn{3}{@{}l}{\code{dstm\_distance\_create\_from\_matrix(ret, n, m)}} \\
\midrule
Requires & \textsc{chk} & $m_{ii} = 0$; $m_{ij} \ge -1$ \\
         & \textsc{res} & $m_{ij} = m_{ji}$ (G2) --- not checked
  \devmark{dev:asymmetric} \\
         & \textsc{res} & $m_{ij} \ne 0$ for $i \ne j$, i.e.\ (G3) --- not
  checked, and $0$ is silently reinterpreted as absence of a channel
  \devmark{dev:zero-distance} \\
Ensures  & & $\cgraph$ is fixed for the session; $\nbr{p}$ follows by
  \cref{def:neighbourhood} and $W$ by \cref{def:path-weight} \\
Blocks   & & no; purely local \\
\bottomrule
\end{tabularx}
\end{table}

\begin{remark}[Choosing the weights]
A channel weight is a modelling decision, not a tuning parameter:
$\lag{p}{q}$ must be a sound lower bound on the simulated latency of every
interaction between the two processes. If anything in the model can cross
faster, the program cannot express it legally without violating
\cref{res:lookahead}. Subject to soundness, larger is better:
\cref{cor:space-drift} shows memory grows with $\diam_p$, but
\cref{thm:pipeline-depth} shows the tolerable drift grows with it too.
\end{remark}

% ===========================================================================
\section{The kernel operations}
\label{sec:api-kernel}

\begin{lstlisting}
int dstm_put(int rank, int region_id, dstm_time_t time, dstm_time_t delta,
             dstm_ptr_t dst, void *src, size_t size);
int dstm_get(int rank, int region_id, dstm_time_t time, dstm_time_t delta,
             void *dst, dstm_ptr_t src, size_t size);
int dstm_get_or_default(int rank, int region_id, dstm_time_t time,
                        dstm_time_t delta, void *dst, dstm_ptr_t src,
                        size_t size, const void *default_value);
\end{lstlisting}

These are the only \glspl{operation} touching the \gls{store}; everything in
\cref{sec:api-darray,sec:api-dgraph} reduces to them.

\begin{table}[htbp]
\centering\small
\begin{tabularx}{\textwidth}{@{}l l >{\raggedright\arraybackslash}X@{}}
\toprule
\multicolumn{3}{@{}l}{\code{dstm\_put(q, q, t, delta, dst, src, size)}
  $\;\equiv\; \Put{p}{q}{t}{\delta}{a}{v}$} \\
\midrule
Requires & \textsc{ast} & \code{rank == region\_id == dst.region}
  (\cref{def:pointer}) \\
         & \textsc{chk} & \code{src} non-\code{NULL}; $\delta \ge 0$;
  if $\delta = 0$ then $q = p$ \\
         & \textsc{res} & $\delta \ge \lag{p}{q}$ (\cref{res:lookahead})
  \devmark{dev:lookahead-uncheck} \\
         & \textsc{res} & $t \ge \clk{p}$ (\cref{res:monotone})
  \devmark{dev:monotone-origin} \\
         & \textsc{res} & no other put targets any cell of
  $\{q\}\times\{t{+}\delta\}\times[a, a{+}\mathit{size})$
  (\cref{res:write-disjoint}) \devmark{dev:write-disjoint} \\
         & \textsc{res} & no get on $(q, t{+}\delta)$ was issued earlier by
  $p$ (\cref{res:phase}) \\
Ensures  & & $\StoreH$ maps the target cells to the bytes of \code{src};
  \code{src} may be reused immediately \\
Frame    & & $\clk{p}$, $\Pend{p}$, $\Net$; and $\Vers{q}$ eventually \\
Blocks   & & no \\
Rules    & & \textsc{Tick}; then \textsc{Put-Local} or \textsc{Put-Remote};
  later \textsc{Apply} \\
\bottomrule
\end{tabularx}
\end{table}

\begin{table}[htbp]
\centering\small
\begin{tabularx}{\textwidth}{@{}l l >{\raggedright\arraybackslash}X@{}}
\toprule
\multicolumn{3}{@{}l}{\code{dstm\_get(q, q, t, delta, dst, src, size)}
  $\;\equiv\; \Get{p}{q}{t}{\delta}{a}{\ell}$} \\
\midrule
Requires & \textsc{ast} & \code{rank == region\_id == src.region} \\
         & \textsc{chk} & \code{dst} non-\code{NULL}; $\delta \le 0$;
  if $\delta = 0$ then $q = p$ \\
         & \textsc{res} & $|\delta| \le \lag{p}{q}$ (\cref{res:staleness})
  \devmark{dev:lookahead-uncheck} \\
         & \textsc{res} & $t \ge \clk{p}$ (\cref{res:monotone}) \\
Ensures  & & \code{dst} holds
  $\StoreH(\cell{q}{t+\delta}{a}) \dots$ (\cref{def:read-result}),
  by \cref{thm:read-soundness} \\
Frame    & & $\clk{p}$, $\Pend{p}$, $\Net$, \code{dst} \\
Blocks   & & \textbf{yes}, until $\Vst{q}{t+\delta} = \Rd$; drives the
  progress engine meanwhile \\
Rules    & & \textsc{Tick}; \textsc{Get-Local}/\textsc{Get-Remote};
  \textsc{Serve-Local}/\textsc{Serve-Remote} \\
\bottomrule
\end{tabularx}
\end{table}

\begin{remark}[Failure modes and how they present]
\label{rem:failure-modes}
The three restrictions fail in visibly different ways, worth knowing when
debugging.
\begin{itemize}[leftmargin=1.4em,itemsep=1pt]
  \item Violating \cref{res:lookahead} or \cref{res:phase} usually reaches
    \rn{Apply-Violation} and \textbf{aborts}, naming the
    \gls{effective-time}. This is the good case: it is loud.
  \item Violating \cref{res:staleness} targets an evicted \gls{version}.
    \rn{Serve} never fires and the program \textbf{hangs} inside
    \code{dstm\_get}, with no diagnostic.
  \item Violating \cref{res:write-disjoint} \textbf{silently} returns
    whichever value arrived last. Nothing detects it, and the result varies
    between runs.
\end{itemize}
\end{remark}

\code{dstm\_get\_or\_default} differs in exactly one respect: if
$t + \delta < 0$ it copies the supplied default and returns without contacting
the owner (\cref{def:get-or-default}). It exists so the first ticks of a
simulation need no special case. Note that on the default path it performs no
bounds or sign checking at all, so a positive $\delta$ paired with a
sufficiently negative $t$ silently succeeds where \code{dstm\_get} would
reject.

% ===========================================================================
\section{Distributed arrays}
\label{sec:api-darray}

\begin{lstlisting}
int dstm_darray_create(dstm_darray_t **ret, size_t element_size,
                       size_t local_elements);
int dstm_darray_create_equal(dstm_darray_t **ret, size_t element_size,
                             size_t total_elements);
int dstm_darray_put(dstm_darray_t *a, size_t i, void *src,
                    dstm_time_t time, dstm_time_t delta);
int dstm_darray_get(dstm_darray_t *a, size_t i, void *dst,
                    dstm_time_t time, dstm_time_t delta);
int dstm_darray_put_range (dstm_darray_t *a, const size_t *idx, size_t n,
                           const void *src, dstm_time_t t, dstm_time_t d);
int dstm_darray_get_range (dstm_darray_t *a, const size_t *idx, size_t n,
                           void *dst, dstm_time_t t, dstm_time_t d);
int dstm_darray_allgatherv(dstm_darray_t *a, void *src, void *dst,
                           dstm_time_t time, const dstm_time_t *deltas);
\end{lstlisting}

A darray adds addressing, not semantics: \code{dstm\_darray\_put(a, g, \dots)}
is \code{dstm\_put} applied to $\mathit{res}(g)$ of \cref{def:darray},
inheriting every restriction of \cref{sec:api-kernel}.

\begin{table}[htbp]
\centering\small
\begin{tabularx}{\textwidth}{@{}l l >{\raggedright\arraybackslash}X@{}}
\toprule
\multicolumn{3}{@{}l}{\code{dstm\_darray\_create(ret, z, m\_p)}} \\
\midrule
Requires & \textsc{chk} & $z \ge 1$, $m_p \ge 1$; a session is open \\
         & \textsc{res} & collective (\cref{res:collective}); no consensus
  guard, so a local allocation failure hangs the job
  \devmark{dev:noncollective} \\
         & \textsc{res} & destroyed collectively, and only once quiescent
  \devmark{dev:destroy-lifetime} \\
Ensures  & & an \gls{allocation} of $z\,m_p$ bytes on each process, plus the
  \gls{partition} of \cref{def:partition} with $b_p$ published by all-gather;
  $\beta_p$ is that allocation's base, chosen locally and exchanged \\
Errors   & & \code{-DSTM\_ENOMEM} iff no free interval of $z\,m_p$ bytes
  exists (\cref{def:alloc-ops}) \\
Blocks   & & yes, two all-gathers \\
\bottomrule
\end{tabularx}
\end{table}

\begin{table}[htbp]
\centering\small
\begin{tabularx}{\textwidth}{@{}l l >{\raggedright\arraybackslash}X@{}}
\toprule
\multicolumn{3}{@{}l}{\code{dstm\_darray\_get/put(a, i, buf, t, delta)}} \\
\midrule
Requires & \textsc{res} & $i < \sum_p m_p$ --- \emph{not checked}; an
  out-of-range index resolves into the last process's region and corrupts
  a neighbouring allocation \devmark{dev:darray-bounds} \\
         & & all restrictions of \code{dstm\_put}/\code{dstm\_get} with
  $q = \mathit{own}(i)$ \\
Ensures  & & operates on the $z$ bytes at $\mathit{res}(i)$
  (\cref{def:darray}) \\
\bottomrule
\end{tabularx}
\end{table}

\subsection*{The collective read--write}

\code{dstm\_darray\_allgatherv} is the one composite \gls{operation}, with an
unusually strict contract: $\mathit{deltas}[p] > 0$ for the caller's own slot
and $\mathit{deltas}[r] < 0$ for every other, zero rejected everywhere.

The reason is \cref{res:phase}. The call publishes the caller's slice then
reads everyone else's. Forcing the local \gls{delta} strictly positive and
remote ones strictly negative places the written and read \glspl{version} at
provably distinct \glspl{effective-time}, so the write-phase/read-phase
discipline cannot be violated by the composite --- \cref{res:phase} enforced
by typing.

\begin{table}[htbp]
\centering\small
\begin{tabularx}{\textwidth}{@{}l l >{\raggedright\arraybackslash}X@{}}
\toprule
\multicolumn{3}{@{}l}{\code{dstm\_darray\_allgatherv(a, src, dst, t, deltas)}} \\
\midrule
Requires & \textsc{chk} & $\mathit{deltas}[p] > 0$;
  $\mathit{deltas}[r] < 0$ for $r \ne p$ \\
         & \textsc{res} & $|\mathit{deltas}|=\worldsize$; identical $t$ on
  every process; $\mathit{deltas}[p] \ge \lag{p}{r}$ and
  $|\mathit{deltas}[r]| \le \lag{p}{r}$ \\
Ensures  & & $\mathit{dst}$ holds every process's slice \emph{except the
  caller's own}, which is never written \\
Blocks   & & yes, once per remote element \\
\bottomrule
\end{tabularx}
\end{table}

That the caller's own slice is left untouched in \code{dst} follows from the
design, not oversight: reading it back would need a local get at the same
\gls{effective-time} as the local put, which \cref{res:phase} forbids.

% ===========================================================================
\section{Distributed graphs}
\label{sec:api-dgraph}

\begin{lstlisting}
int dstm_dgraph_create(dstm_dgraph_t **ret, size_t vertex_state_size,
                       size_t edge_state_size, size_t local_vertices,
                       const size_t *local_out_degrees,
                       const size_t *local_edge_dsts,
                       const void *local_vertex_states,
                       const void *local_edge_states);
int dstm_dgraph_vertex_get(dstm_dgraph_t *g, size_t v, void *dst,
                           dstm_time_t time, dstm_time_t delta);
int dstm_dgraph_edge_put  (dstm_dgraph_t *g, size_t u, size_t v, void *src,
                           dstm_time_t time, dstm_time_t delta);
int dstm_dgraph_neighbors_get(dstm_dgraph_t *g, size_t u,
                              size_t *out, size_t *io_count);
int dstm_dgraph_derive_distance(dstm_dgraph_t *g, dstm_distance_t **ret,
                                int default_delta);
\end{lstlisting}

A \glslink{dgraph}{dgraph} holds \emph{two} kinds of state, and keeping them apart is the key to
using the API.

\begin{description}[leftmargin=1.3cm,style=nextline,font=\normalfont\bfseries]
\item[Topology] --- the vertex set, the edge set and the adjacency structure.
  Frozen at construction, held by every neighbouring \gls{process} as a
  \gls{topology-replica}, and \emph{not versioned}. Queries such as \code{neighbors\_get} and
  \code{degree\_get} therefore carry no $(t, \delta)$ and involve no
  communication and no blocking.
\item[State] --- the per-vertex and per-edge payloads. Versioned exactly like
  darray elements, with identical restrictions.
\end{description}

\begin{remark}[Edges are owned by their source]\label{rem:edge-ownership}
The state of edge $(u,v)$ lives on the owner of $u$, which makes the natural
idiom legal: reading edge $(u,v)$ with $u$ local is a local get at $\delta =
0$, writing edge $(v,u)$ with $v$ remote is a remote put at $\delta =
+\lag{}{}$. Reading $(x,y)$ as ``the inbox at $x$ for traffic from $y$'' puts
every \gls{operation} on the legal side of the ownership boundary and gives
each \gls{cell} one writer, discharging \cref{res:write-disjoint} structurally
--- the \glslink{mailbox}{mailbox pattern} in graph form.
\end{remark}

\begin{table}[htbp]
\centering\small
\begin{tabularx}{\textwidth}{@{}l l >{\raggedright\arraybackslash}X@{}}
\toprule
\multicolumn{3}{@{}l}{\code{dstm\_dgraph\_edge\_get/put(g, u, v, buf, t, delta)}} \\
\midrule
Requires & \textsc{chk} & $u, v <$ total vertices; edge $(u,v)$ declared,
  else \code{-DSTM\_ENOTFOUND} \\
         & \textsc{chk} & the owner of $u$ is the local process or a topology
  neighbour, else \code{-DSTM\_ENOTLOCAL} \\
         & \textsc{res} & the restrictions of \code{dstm\_put}/\code{dstm\_get}
  with $q = \mathit{own}(u)$ \\
Ensures  & & operates on the edge-state slot of $(u,v)$ \\
Blocks   & & get: yes. put: no. Resolution of $(u,v)$ itself: never \\
\bottomrule
\end{tabularx}
\end{table}

\code{-DSTM\_ENOTLOCAL} is a genuine partiality, not an error: a vertex owned
by a non-neighbouring process is a valid global identifier whose adjacency
this process does not replicate.

\begin{table}[htbp]
\centering\small
\begin{tabularx}{\textwidth}{@{}l l >{\raggedright\arraybackslash}X@{}}
\toprule
\multicolumn{3}{@{}l}{\code{dstm\_dgraph\_derive\_distance(g, ret, d)}} \\
\midrule
Requires & \textsc{chk} & $d > 0$; collective \\
Ensures  & & $\Dmat_{ii} = 0$; $\Dmat_{ij} = d$ if some edge joins a vertex of
  $i$ to one of $j$ \emph{in either direction}; no channel otherwise \\
         & & the result is a well-formed \gls{comm-graph}
  (\cref{prop:derive-wf}) \\
Blocks   & & yes, one all-gather \\
\bottomrule
\end{tabularx}
\end{table}

\begin{proposition}[Derived matrices are well formed]
\label{prop:derive-wf}
The graph encoded by \code{dstm\_dgraph\_derive\_distance} satisfies
(G1)--(G3).
\end{proposition}

\begin{proof}
(G1) holds because the diagonal is written unconditionally to $0$,
independently of any edge data, so intra-process edges cannot perturb it.
(G2) holds because the adjacency predicate is the symmetric closure
$\mathit{touch}(i,j) \vee \mathit{touch}(j,i)$ of the ``some edge crosses''
relation, and the assigned weight is the same $d$ on both sides. (G3) holds
because $d > 0$ is checked and every channel weighs exactly $d$.

\end{proof}

\Cref{prop:derive-wf} is why the dgraph path has never exhibited
\cref{dev:zero-distance} or \cref{dev:asymmetric}: it cannot construct a graph
that triggers them. Only the hand-written \code{create\_from\_matrix} path
can.

% ===========================================================================
\section{Legal idioms}
\label{sec:api-idioms}

The restrictions are easiest to discharge by construction. Three patterns do
so.

\subsection*{Push}

\begin{lstlisting}
for (t = 0; t < T; t++) {
        /* consume: own inbox at the current tick */
        for (s in neighbours) dstm_darray_get(a, cell(me, s), &in[s], t, 0);
        compute();
        /* produce: neighbour inboxes, one lag ahead      */
        for (d in neighbours) dstm_darray_put(a, cell(d, me), &out[d], t, D(me,d));
}
\end{lstlisting}

Discharges \cref{res:lookahead} ($\delta = \lag{}{}$ exactly),
\cref{res:staleness} (local reads have $\delta = 0$),
\cref{res:write-disjoint} (cell $(d, \mathit{me})$ has $\mathit{me}$ as its
only writer) and \cref{res:phase} (all \glspl{get} precede all \glspl{put} at each tick, and
they target different \glspl{effective-time} anyway). Consumer tick is
$t + \lag{}{}$ by \cref{thm:pipeline-depth}. One message per edge per tick.

\subsection*{Pull}

\begin{lstlisting}
for (t = 0; t < T; t++) {
        /* produce: own region at the current tick */
        dstm_darray_put(a, cell(me, me), &out, t, 0);
        /* consume: neighbours' regions, one lag behind  */
        for (s in neighbours) dstm_darray_get(a, cell(s, s), &in[s], t, -D(me,s));
        compute();
}
\end{lstlisting}

Same discharge, with the put and the get exchanging roles. Note the put must
precede the gets in \gls{program-order} at each tick: they are at different
effective times so \cref{res:phase} is not strictly at issue, but a \gls{process} that
reads before publishing stalls its neighbours for a full round trip.

\subsection*{Mailbox indexing}

Both idioms index the shared structure by \emph{(destination, source)}. This
is what makes \cref{res:write-disjoint} structural rather than a proof
restriction: cell $(d,s)$ is written only by $s$ and read only by $d$, so no
two writers can ever meet. Any application that can be phrased this way gets
determinism for free, and this is the recommended default.

\begin{remark}[Publish every tick, including empties]\label{rem:publish-empty}
Under \cref{conv:no-inherit} an unwritten \gls{version} still exists,
zero-filled, and a reader will observe it. The subtle failure is the opposite:
a producer skipping a tick because it has nothing to send leaves its consumer
reading zeroes, usually the intended meaning of ``nothing arrived''.
Publishing every outbound direction every tick also guarantees one writer per
edge version, the cleanest way to satisfy \cref{res:write-disjoint}.
\end{remark}